\section{Bayesian Receiver Design}
\label{sec:algorithm}

The receiver consists of a linear estimator and a channel and activity
denoiser. The linear estimator combines the received pilots with the
messages from the denoiser, while the denoiser exploits the channel
prior and device activity distribution to update the effective
channels. Following the message-passing structure in
\cite{cai2026}, the two modules exchange Gaussian extrinsic messages.
We unfold these updates into $L$ layers and introduce trainable
parameters in each layer \cite{wei2026}. The message calculations
are described below, with $J=KN$ and $C=MN$.

\subsection{Linear Estimation}

At layer $\ell$, the input message for each column of $\mathbf X$
has mean $\mathbf M_\ell$ and covariance $w_\ell\I_J$.
Combining this message with the likelihood in
\eqref{eq:working-posterior} gives a Gaussian distribution.
By completing the square, its mean and covariance are obtained as
\begin{align}
 \mathbf M_{{\rm L},\ell}&=\mathbf M_\ell+\bm\Delta_\ell,
 \label{eq:twc-linear-mean}\\
 \bm\Delta_\ell&=w_\ell\mathbf A_{\rm c}^{\mathsf H}
 (\sigma_{\rm c}^2\I_Q+w_\ell\mathbf A_{\rm c}\mathbf A_{\rm c}^{\mathsf H})^{-1}
 (\mathbf Y-\mathbf A_{\rm c}\mathbf M_\ell),
 \label{eq:twc-linear-correction}\\
 \mathbf C_\ell
 &=\left(w_\ell^{-1}\I_J+
 \frac{\mathbf A_{\rm c}^{\mathsf H}\mathbf A_{\rm c}}{\sigma_{\rm c}^2}
 \right)^{-1}.
 \label{eq:twc-linear-covariance}
\end{align}
We approximate the covariance by $t_\ell\I_J$, where
$t_\ell=\tr(\mathbf C_\ell)/J$. The extrinsic message is calculated
by dividing this Gaussian approximation by the input message
\cite{minka2001}. Its variance and mean are respectively given by
\begin{equation}
 v_\ell^{(0)}=(t_\ell^{-1}-w_\ell^{-1})^{-1},\qquad
 \mathbf R_\ell^{(0)}=v_\ell^{(0)}
 \left(\frac{\mathbf M_{{\rm L},\ell}}{t_\ell}
       -\frac{\mathbf M_\ell}{w_\ell}\right).
 \label{eq:twc-linear-extrinsic}
\end{equation}
Letting $f_\ell=t_\ell/w_\ell$ and $d_\ell=1-f_\ell$, we can
equivalently write
\begin{equation}
 v_\ell^{(0)}=w_\ell f_\ell/d_\ell,\qquad
 \mathbf R_\ell^{(0)}=\mathbf M_\ell+\bm\Delta_\ell/d_\ell.
 \label{eq:twc-query-uncalibrated}
\end{equation}
To adjust the messages during receiver training, introduce a positive
calibration factor $\gamma_\ell=e^{c_{\gamma,\ell}}$ and set
\begin{equation}
 \xi_\ell=\gamma_\ell f_\ell/d_\ell,\qquad
 v_\ell=w_\ell\xi_\ell,\qquad
 \mathbf R_\ell=\mathbf M_\ell+(1+\xi_\ell)\bm\Delta_\ell.
 \label{eq:twc-query-calibrated}
\end{equation}
The mean $\mathbf R_\ell$ and variance $v_\ell$ form the input to the
channel and activity denoiser. When $\gamma_\ell=1$,
\eqref{eq:twc-query-calibrated} reduces to
\eqref{eq:twc-query-uncalibrated}.

\subsection{Channel and Activity Denoising}

The denoiser estimates the channel and activity of each terminal,
and combines them to obtain its effective channel.
For a user block $r\in\mathbb C^D$ of $\mathbf R_\ell$, the Gaussian
input message is proportional to $\CN(r;\alpha h,v\I_D)$.
Combining it with the channel and activity priors gives
\begin{equation}
 q_0(r;v)=\CN(r;0,v\I_D),\qquad
 q_1(r;v)=\int\CN(r;h,v\I_D)q_{\rm tr}(dh),
 \label{eq:twc-local-densities}
\end{equation}
where $q_0$ and $q_1$ correspond to the inactive and active
hypotheses, respectively. The activity probability and active-channel
conditional mean are
\begin{equation}
 p^\star(r,v)=\frac{\rho q_1(r;v)}
                   {(1-\rho)q_0(r;v)+\rho q_1(r;v)},\qquad
 \mu^\star(r,v)=\frac{\int h\,\CN(r;h,v\I_D)q_{\rm tr}(dh)}{q_1(r;v)}.
 \label{eq:twc-local-targets}
\end{equation}
Since the effective channel is zero under the inactive hypothesis,
its posterior mean is given by \cite[eq.~(35)]{cai2026}
\begin{equation}
 \E[x\mid r]=p^\star(r,v)\mu^\star(r,v).
 \label{eq:twc-mixture-mean}
\end{equation}
The integrals in \eqref{eq:twc-local-targets} involve the distribution
of the complete multiantenna channel. We approximate these two outputs
using a score-based channel denoiser and an activity classifier.

Let $u_{k\ell}\in\mathbb R^{2D}$ denote the real and imaginary parts
of the $k$th user query. Based on the score representation developed
in Section~\ref{sec:learning}, the channel estimate is
\begin{equation}
 \mu_{k\ell}=u_{k\ell}
       +a_\ell\frac{v_\ell}{2}s_{\theta_\ell}(u_{k\ell},v_\ell),
 \qquad a_\ell=e^{c_{\alpha,\ell}}.
 \label{eq:twc-conditional-mean}
\end{equation}
Here $s_{\theta_\ell}$ is the score network in layer $\ell$,
$v_\ell/2$ is the noise variance per real coordinate, and $a_\ell$
controls the denoising correction. The activity classifier uses the
query and channel estimate to compute
\begin{equation}
 b_{k\ell}=\logit(\rho)+
 g_{\phi_\ell}(u_{k\ell},v_\ell,\mu_{k\ell},\rho),\qquad
 p_{k\ell}=\sigmoid(b_{k\ell}),
 \label{eq:twc-activity}
\end{equation}
where $g_{\phi_\ell}$ denotes the learned correction to the prior
log odds. Substituting these estimates into
\eqref{eq:twc-mixture-mean} gives
\begin{equation}
 \widehat x_{k\ell}=p_{k\ell}\mu_{k\ell}.
 \label{eq:twc-estimate}
\end{equation}
The user estimates are rearranged as $\widehat{\mathbf X}_\ell$
according to \eqref{eq:mmv-indexing}. Within each layer, all users
share the score network and activity classifier.

\subsection{Extrinsic Message Update}

To provide feedback to the linear estimator, we assign a shared
variance $\tau_\ell$ to the Gaussian output message with mean
$\widehat{\mathbf X}_\ell$. Following the residual-based construction
in \cite[eq.~(12)]{wei2026}, we first compute
\begin{equation}
 e_\ell=\frac{\|\mathbf Y-\mathbf A_{\rm c}\widehat{\mathbf X}_\ell\|_F^2
                   -QC\sigma_{\rm c}^2}
                  {C\|\mathbf A_{\rm c}\|_F^2}.
 \label{eq:twc-residual-scale}
\end{equation}
The residual is formed from the activity-weighted channel estimates.
We obtain the feedback scale $\tau_\ell$ by calibrating $e_\ell$
and constraining it to $0<\tau_\ell<v_\ell$, as specified in
Section~\ref{sec:twc-damping}. This provides a positive outgoing
Gaussian precision.

Dividing the Gaussian output by the input message
$\CN(\mathbf X;\mathbf R_\ell,v_\ell\I)$ gives
\begin{equation}
 \widetilde{\mathbf M}_{\ell+1}
 =\frac{\widehat{\mathbf X}_\ell-q_\ell\mathbf R_\ell}{1-q_\ell},
 \qquad
 \widetilde w_{\ell+1}=\frac{v_\ell q_\ell}{1-q_\ell},
 \qquad q_\ell=\tau_\ell/v_\ell.
 \label{eq:twc-feedback-proposal}
\end{equation}
The resulting message is damped before it becomes the input to the
next layer. Thus the activity and channel estimates jointly determine
both the observation residual and the information fed back to the
linear estimator.
